
\section{Tense, Aspect and Modality (TAM)}

\begin{frame}{TAM}
\begin{itemize}
\item  We need to distinguish grammatical expression from 
  meaning
  \begin{itemize}
  \item  Tense vs Time
  \item  Grammatical Aspect vs Semantic Aspect
  \item  Mood vs Modality
  \item  Surface Case vs Deep Case
  \end{itemize}
\item  The relation between them is refered to as
  \begin{itemize}
  \item  linking; syntax-semantics interface; grammar
  \end{itemize}
\end{itemize} 
 

\end{frame}

\begin{frame}{How Universal is Tense?}
\begin{itemize}
\item  Grammatical tense is different from semantic time
\item  English has \iz{past/non-past}
\item  Latin marks \iz{past/present/future}
\item  Chibemba (Bantu) has \txx{metrical tense}
  \begin{multicols}{2}
  \begin{itemize}
  \item  Remote Past ($<$ yesterday)
  \item  Removed Past (yesterday)
  \item  Near Past (today)
  \item  Immediate Past (past few hours)
  \item Immediate Future (next few hours)
  \item Near Future (today)
  \item Removed Future (tomorrow)
  \item Remote Future ($>$ tomorrow)
  \end{itemize}
\end{multicols}
\end{itemize}

\end{frame}

\begin{frame}{Tense and Time}
\begin{itemize}
\item  Locate a situation to with respect to a point in time
  \begin{itemize}
  \item  S = speech point
  \item  R = reference time
  \item  E = event time
  \end{itemize}
\item   Hans Reichenbach (1947)
\item Simple Tense
\begin{itemize}
\item  Past ($R = E < S$) \eng{saw} \hfill
\begin{tabular}[t]{ccc|ccc|ccc}
  \mc{3}{c}{past} &  \mc{3}{c}{present} & \mc{3}{c}{future} \\
&R=E&&&S&&&& \\ \hline
&&&&&&& \\ 
\end{tabular}
\item  Present ($R = S = E$) \eng{see} \hfill
\begin{tabular}[t]{ccc|ccc|ccc}
  \mc{3}{c}{past} &  \mc{3}{c}{present} & \mc{3}{c}{future} \\
&&&&S=R=E&&&&  \\ \hline
&&&&&&& \\ 
\end{tabular}
\item  Future ($S < R = E$) \eng{will see}\hfill
\begin{tabular}{ccc|ccc|ccc}
  \mc{3}{c}{past} &  \mc{3}{c}{present} & \mc{3}{c}{future} \\
&&&&S&&&R=E& \\ \hline
&&&&&&& \\ 
\end{tabular}
\end{itemize}
\end{itemize}
\end{frame}

\begin{frame}{Complex Tense}

\begin{itemize}
\item  Past Perfect ($E < R < S$) \eng{had seen} \hfill
\begin{tabular}[t]{ccc|ccc|ccc}
  \mc{3}{c}{past} &  \mc{3}{c}{present} & \mc{3}{c}{future} \\
E&R&&&S&&&& \\ \hline
&&&&&&& \\ 
\end{tabular}
\\ \eng{By 1939 my Father had seen many arrests}
\item  Future Perfect ($S< E < R$) \eng{will have seen} \hfill
\begin{tabular}[t]{ccc|ccc|ccc}
  \mc{3}{c}{past} &  \mc{3}{c}{present} & \mc{3}{c}{future} \\
&&&&S&&&E&R \\ \hline
&&&&&&& \\ 
\end{tabular}
\\ \eng{By 2039 my son will have seen many things}
\end{itemize}

\end{frame}

\begin{frame}{Aspect in English}
\begin{itemize}
\item  Finer grained talking about time!
\item  \txx{Progressive}  is used for ongoing processes (unfinished)
  \begin{itemize}
  \item  \txx{Past Progressive} \eng{I was building the building}
  \item  \txx{Present Progressive} \eng{I am building the building}
  \item  \txx{Future Progressive} \eng{I will be building the building}  %(E < R = S) ???
  \end{itemize}
\item  \txx{Perfect} compares the time to the reference point
 \begin{itemize}
  \item  \txx{Past Perfect} \eng{I had built the building} ($E<R<S$)
  \item  \txx{Present Perfect} \eng{I have built the building}  ($E<R=S$)
  \item  \txx{Future Perfect} \eng{I will have built the building}  ($S<E<R$)
  \end{itemize}
\end{itemize}

\end{frame}

\begin{frame}{Aspect more Generally}

\begin{itemize}
\item  \txx{Perfective} focuses on the end point
  \begin{itemize}
  \item  \txx{Completive} \eng{I built the building}
  \item  \txx{Experiential} \eng{I have built the building} %(E < R = S) ???
  \end{itemize}
\item  \txx{Imperfective} 
  \begin{itemize}
  \item  \txx{Progressive} \eng{I was listening/I am listening}
  \item  \txx{Habitual} \eng{I listen to the Goon Show}
\end{itemize}
\item Different languages grammaticalize different things
\end{itemize}

% \item  Progressive
% \item  Habitual
% \item  Completive
% \item  Experiential

\end{frame}

\begin{frame}{Mood and Modality}
\begin{itemize}
\item  Modality expresses varying degrees of the speaker's 
commitment and belief
\item In English it is typically expressed by an auxiliary verb.
\begin{exe}
  \ex \eng{She \ul{has} left by now.}
  \ex \eng{She \ul{must} have left by now.}
  \ex \eng{She \ul{could} have left by now.}
  \ex \eng{She \ul{needn't} have left by now.}
  \ex \eng{She \ul{couldn't} have left by now.}
  \ex \eng{She \ul{has} to leave by now.}
  \ex \eng{She \ul{must} leave by now.}
  \ex \eng{She \ul{can} leave now. }
\end{exe}
\end{itemize}

\end{frame}

\begin{frame}{Other means of expression}
\begin{itemize}
\item Explicit External Verb
  \begin{exe}
    \ex \eng{I \ul{know} that $S$}
    \ex \eng{I \ul{believe} that $S$}
  \end{exe}
\item Adverb or Adjective
  \begin{exe}
    \ex \eng{It is \ul{certain that} $S$}
    \ex \eng{It is \ul{likely that} $S$}
    \ex \eng{I will \ul{probably} $S$}
    \ex \eng{I will \ul{definitely} $S$}
  \end{exe}

\end{itemize}

\end{frame}

\begin{frame}{Knowledge vs Obligation}
\begin{itemize}
\item  \txx{Epistemic modality}: Speaker signals degree of  
knowledge.
\begin{exe}
  \ex \eng{You can drive this car} \cmm{It is possible}
\end{exe}
\item  \txx{Deontic modality}: Speaker signals his/her attitude  
to social factors of obligation and permission.
\begin{itemize}
\item \txx{Permission}
  \begin{exe}
    \ex \eng{You can drive this car}  \cmm{You have permission to}
    \ex \eng{You may drive this car}  
  \end{exe}
\item \txx{Obligation}
  \begin{exe}
    \ex \eng{You must drive this car}  \cmm{You have an obligation to}
    \ex \eng{You ought to drive this car} 
  \end{exe}
% \item  Question: How do you use `must', `confirm', and 
% `sure' in Singlish? ???
\end{itemize}
\item [?] Find examples of each type (deontic/epistemic), expressed in different ways (auxiliary, verb, adverb/adjective) in \Story{}\task
\end{itemize}

\end{frame}

\begin{frame}{Other modalities}

There are other possible modalities \citep[p290 (4)]{Kroeger:2022}:
\begin{itemize}
\item \txx{bouletic} concerns what is possible or necessary,
given a person's desires
\item \txx{dynamic} concerns properties and dispositions of persons (such as ability)
\item \txx{teleogical} concerns achieving a goal 
\end{itemize}

\begin{exe}
\item \eng{It has to be raining.}
   [after observing people coming inside with wet umbrellas] \txx{epistemic}
 \item \eng{Visitors have to leave by six pm}.
   [hospital regulations] \txx{deontic}
 \item \eng{John has to work hard if he wants to retire at age 50.}
   [to attain desires] \txx{bouletic}
 \item \eng{I have to sneeze.}
   [given the current state of one's nose] \txx{dynamic}
 \item \eng{To get home in time, you have to take a taxi.}
   [in order to achieve the stated purpose] \txx{teleological}
\end{exe}
\end{frame}

\begin{frame}{Possible Worlds}

\begin{itemize}
\item We can analyze these in terms of \txx{possible worlds}

\item We mark how close a hypothetical case is to reality:
  \begin{exe}
    \ex \eng{It must be/might be/is/can't be \ul{hot outside}}
  \end{exe}
  And model it as the degree of overlap of the worlds
\item Similarly for \txx{conditionals} (condition/consequence)
  \begin{exe}
    \ex \eng{If it is Singapore, it will be hot outside}
    \ex \eng{If it were Singapore, it would be hot outside}
    \ex \eng{If you should go to Singapore, take some cool clothes}
  \end{exe}
\end{itemize}

% \end{frame}

\begin{frame}{Real vs Hypothetical}
% \begin{itemize}
% \item  \txx{Realis} is used for things that occur
% \item  \txx{Irrealis} is used for things that are not claimed to 
% occur (hypotheticals, negation, future)
% \item  English doesn't mark this normally
%   \begin{exe}
%     \ex \eng{If I were to go} \textnormal (subjunctive)
%   \end{exe}
% \item  What about Singlish? ???
%  \begin{exe}
%     \ex \eng{I got go.}
%     \ex \eng{I sure confirm go.}
%     \ex \eng{I maybe go.}
%   \end{exe}
% \end{itemize}

\end{frame}

\begin{frame}{Mood more Generally}
\MyLogo{See \citet[Chapter 5]{Saeed:2015} for more examples}
\begin{itemize}
\item Grammatical Inflection used to mark modality is called \txx{mood}
  \begin{itemize}
  \item \txx{indicative} expresses factual statements
  \item \txx{conditional} expresses events dependent on a condition
  \item \txx{imperative} expresses commands
  \item \txx{injunctive} expresses pleading, insistence, imploring
  \item \txx{optative} expresses hopes, wishes or commands 
  \item \txx{potential} expresses something likely to happen
  \item \txx{subjunctive} expresses  hypothetical events; opinions or emotions
  \item \txx{interrogative} expresses questions
\end{itemize}
\item In English most things are marked syntactically: \vspace*{-1ex}
    \begin{exe}
      \ex  \eng{I \ul{am} good}  \hfill indicative
      \ex  \eng{\ul{Am} I good?} \hfill interrogative
    \ex \eng{\ul{Be} good!} \hfill imperative
    \ex \eng{If I \ul{were} a rich man} \hfill subjunctive
  \end{exe}
\item [?] Find examples of non-indicative sentences in \Story{}\task
\end{itemize}

\end{frame}

\begin{frame}{Evidentiality}
\MyLogo{Skip if too long}
\newcommand{\TIPA}[1]{\textipa{\mtcitestyle{#1}}}

\begin{itemize}
\item  Some languages must show you gained the evidence
\\item  \txx{nonvisual sensory}:  speaker felt the sensation
  \begin{itemize}
  \item \ipa{/pʰa·békʰ-\ul{ink’e/}} ``burned, I felt it''
  \end{itemize}
\item  \txx{inferential}: speaker saw circumstantial evidence 
  \begin{itemize}
  \item  \ipa{/pʰa·bék-\ul{ine}/}  ``must have burned''
  \end{itemize}
\item  \txx{hearsay (reportative)}:   speaker is reporting what was told
  \begin{itemize}
  \item  \ipa{/pʰa·békʰ-\ul{·le}/} ``burned, they say''
  \end{itemize}
\item  \txx{direct knowledge}:   speaker has direct evidence, probably visual 
  \begin{itemize}
  \item \ipa{/pʰa·bék-\ul{a}/} ``burned, I saw it''
  \end{itemize}
\item Examples from Eastern Pomo (McLendon 2003)
\end{itemize}

\end{frame}

\begin{frame}{Evidentiality in English}

We can, and often do, mark evidentiality in English, although it is
not strongly grammaticalized.

\begin{exe}
\ex \eng{Bob is hungry.}
\ex \eng{Bob looks hungry.}
\ex \eng{Bob seems hungry.}
\ex \eng{Bob is apparently hungry.}
\ex \eng{Bob would be hungry by now.}
\ex \eng{Look at those clouds! It's going to rain!}
\ex \eng{Look at those clouds! $^\#$ It will rain!.}
\end{exe}


\end{frame}

